% TLDR: Explain what additional data buys and how dependence constrains precision claims.
\section{Precision, Selection, and Decision Rules}
\label{sec:inference}
The three dimensions of the study address different uncertainties.
More distinct tasks cover a broader set of inputs and increase the precision of task-averaged comparisons.
More solver configurations probe different output distributions.
More distinct raters probe variation between raters, whereas more ratings of one artifact sharpen information about that artifact's rating distribution.
The product of these counts is not the number of independent observations.

\subsection{A transparent planning calculation}
As an illustration, suppose the human-specific contribution to normalized loss has standard deviation $0.10$.
For independent, representative raters, the normal-approximation 95\% half-width of that contribution would be
\begin{equation}
 1.96\frac{0.10}{\sqrt H}
 =\begin{cases}0.080,&H=6,\\0.049,&H=16.\end{cases}
\end{equation}
This is a hypothetical component of uncertainty, not an observed total interval.
If repeated ratings share a human-specific effect, thousands of ratings by the same few raters do not remove it.
The calculation motivates recruiting more distinct raters, while the actual allocation must be checked against pilot variance and workload.

Increasing ratings per artifact from three to eight reduces an independent mean's standard-error factor by $\sqrt{3/8}\approx0.61$.
That is a 39\% reduction under the stated independence and common-variance approximation, not a guaranteed 39\% reduction in the final study interval.
Balanced reuse of raters creates cross-artifact dependence, and an exact finite-panel analysis differs from independent population sampling.
These conditions are why the dense subset supplements broad task coverage instead of replacing it.

For correlations, Fisher's approximation for independent bivariate-normal pairs is
\begin{equation}
 \operatorname{SE}(\operatorname{atanh}\widehat\rho)\approx1/\sqrt{n-3}.
\end{equation}
In the actual study, multiple outputs share a task and ratings share a rater; substituting $n=14{,}848$ into that formula is invalid.
Nor can additional data be expected to raise the true correlation: it improves estimation of whatever association exists.
We therefore size the study against paired proper-loss precision and expert variation, and report correlation intervals under a dependence-aware analysis.

\subsection{Estimating distribution error}
The raw RPS includes the disagreement term in Equation~\ref{eq:decomposition}.
For a fixed artifact and threshold, let $X_h=\mathbf1\{R_h\leq c\}$ and let $s_X^2$ be the sample variance with denominator $m-1$.
Under independent draws from the specified expert population,
\begin{equation}
 \widehat D_c=(Q_c-\overline X)^2-\frac{s_X^2}{m}
 \label{eq:corrected}
\end{equation}
is unbiased for $(Q_c-P_c)^2$.
The subtraction removes the variance of the empirical human cumulative probability.
We average across thresholds and then artifacts.
Individual corrected estimates can be negative and are not clipped before aggregation.
Their intervals, rather than the sign of an individual estimate, determine what distribution error remains compatible with the data.

If the estimand instead concerns a fixed panel of $H$ raters sampled uniformly without replacement, the subtraction is $(1-m/H)s_X^2/m$.
That correction targets the finite-panel mean, not the wider population automatically.
More complicated assignment and nonresponse require corresponding design-based or model-based adjustments.
The dense subset makes this diagnostic more informative but cannot remove uncertainty about unrepresented experts.

\subsection{Uncertainty and comparison}
Crossed task and rater identities require inference that respects both sources of dependence.
Methods for resampling crossed arrays motivate independent reweighting of identities under appropriate assumptions \citep{owen2012bootstrapping}.
The exact sparse-assignment procedure, coverage checks, and treatment of small rater counts must be specified using pilot data before the test is inspected.
We report fixed-panel analyses separately from analyses extending to an expert population.
Task resampling estimates variation over tasks like those in the benchmark; for the complete fixed benchmark it is not uncertainty caused by missing benchmark task identities.
The four chosen solver configurations remain fixed in the primary analysis.

Primary comparisons use paired, artifact-weighted RPS differences between a constant training-estimated distribution, a fixed rubric/reference-assisted judge, a correlation-selected judge, and an RPS-selected judge.
We also include a non-LLM theorem-count baseline to test whether a semantic model improves on a simple structural cue.
Let $v=\min(1,N_{\mathrm{candidate}}/N_{\mathrm{reference}})$ count theorem statements without inspecting their proofs.
For $u=5v$, place probability $1-(u-\lfloor u\rfloor)$ at $\lfloor u\rfloor$ and the remaining mass at the next category, with all mass at five when $u=5$.
This defines an ordinal probability forecast for comparison; it does not assert that counts establish coverage.
A reference with no countable statements is flagged as an invalid baseline input rather than silently assigned a score.
The declaration-counting rule is fixed in the pilot.
Each probabilistic condition must produce the same six-category output format.
Converting a historical scalar judge into a probability forecast is a new baseline condition and must be documented.
Matched artifacts and rating assignments reduce unnecessary comparison noise.
Unequal model calls, tokens, and runtime must be reported rather than hidden as equivalent budgets.

The constant forecast tests whether input-dependent predictions add information beyond the training rating frequencies.
The count baseline tests whether a cheap structural signal explains the apparent gain.
The fixed-rubric condition tests the benefit of fitting or selection, while the correlation-selected condition tests the choice of selection objective.
For the last contrast, both selection procedures must search the same predefined candidate set with the same information access and computational budget.
Otherwise, an observed difference could reflect search resources or a different model family rather than the objective.

The pilot must set a minimum meaningful improvement, a tolerable distribution error, repeat-instability and failure limits, and the corresponding interval rule.
If a claim requires success on all four configurations, inference must account for those four simultaneous claims.
These thresholds are currently unset; a statistically nonzero improvement alone is not the proposed definition of practical success.
Later changes after test inspection require fresh evidence or an explicitly exploratory label.
